\subsection{Actual finite affine maxima and derived outgoing slope inequality}

These are general finite real arithmetic theorems. Their inputs are original
finite families of real numbers and slopes, not a harmonic function, local
profile or slope-balancing certificate. The analytic consumers must still
prove their OWN genuine affine-ray formulas and individual balancing.

\begin{theorem}[Actual affine maxima have a genuine common positive active neighborhood]
\label{analysis:actual-finite-affine-active-max}
\lean{EnvelopeProperty.FiniteAffineMax.active,
EnvelopeProperty.FiniteAffineMax.active_nonempty,
EnvelopeProperty.FiniteAffineMax.activeSlope,
EnvelopeProperty.FiniteAffineMax.exists_pos_eventual_max}
Let $I$ be an ACTUAL finite nonempty set, $D$ any ACTUAL finite set,
$b_i\in\mathbb R$, and $s_{i,d}\in\mathbb R$. Put
$B=\max_{i\in I}b_i$, $I_0=\{i\in I:b_i=B\}$ and
$M_d=\max_{i\in I_0}s_{i,d}$. The actual active set $I_0$ is nonempty.
There EXISTS an actual $\varepsilon>0$ such that for EVERY $d\in D$ and
EVERY $0\leq t<\varepsilon$,
\[
 \max_{i\in I}(b_i+s_{i,d}t)=B+M_dt.
\]
No common neighborhood, active-profile formula or bound on slopes is assumed.
The empty $D$ case is included.
\leanok
\end{theorem}
\begin{proof}
The true maximum of a finite nonempty set is attained, so choose an ACTUAL
$k\in I_0$. For every original inactive $i\in I\setminus I_0$ and
$d\in D$, the actual affine difference
$(b_i+s_{i,d}t)-(b_k+s_{k,d}t)$ is strictly negative at $t=0$.
Its genuine real continuity gives a neighborhood of zero where it stays
strictly negative. There are only finitely many original pairs $(i,d)$,
so their actual finite intersection is still a neighborhood of zero.
Choose a genuine positive metric-ball radius $\varepsilon$ inside it.
For $0\leq t<\varepsilon$, EVERY inactive affine term is smaller than the
actual active $k$ term, and thus cannot increase the maximum.
For active $i$ the actual intercept is $B$; because $t\geq0$, its term is
at most $B+M_dt$. The true finite maximum defining $M_d$ is attained by
some actual active index, giving equality of the whole original maximum.
If there are no inactive pairs the same neighborhood argument uses the
whole real space, still providing an actual positive radius.
\end{proof}

\begin{theorem}[Individual actual balancing gives nonnegative total active-max slope]
\label{analysis:actual-active-max-slope-sum}
\lean{EnvelopeProperty.FiniteAffineMax.sum_activeSlope_nonneg}
\uses{analysis:actual-finite-affine-active-max}
With the SAME original finite families, suppose EVERY original member has
actual zero total outgoing slope: $\sum_{d\in D}s_{i,d}=0$ for EVERY $i\in I$.
Then $\sum_{d\in D}M_d\geq0$.
Together with the preceding genuine neighborhood formula this proves the
outgoing slope inequality for the ACTUAL maximum of these original affine
profiles; it does not assert a general analytic maximum principle.
\leanok
\end{theorem}
\begin{proof}
Choose a genuine original active $k\in I_0$. For every $d\in D$, the true
finite maximum gives $s_{k,d}\leq M_d$. Sum these actual inequalities over
the original finite direction set, and use the GIVEN individual equality
$\sum_ds_{k,d}=0$ to get the claimed nonnegative sum. Here individual
balancing is an explicit arithmetic hypothesis: analytic consumers must
derive it for their original functions, as in the physical rational-unit
branch-balancing theorem. No inferred geometry is hidden in this lemma.
\end{proof}
